\documentclass[11pt,letterpaper]{report}
\usepackage[utf8]{inputenc}
\usepackage[T1]{fontenc}
\usepackage[spanish,es-tabla]{babel}
\usepackage[letterpaper,left=27mm,right=23mm,top=32mm,bottom=25mm,headheight=38pt,headsep=12pt]{geometry}
\usepackage[scaled]{helvet}
\renewcommand{\familydefault}{\sfdefault}
\usepackage{amsmath,amssymb,graphicx,booktabs,longtable,tabularx,array,enumitem,xcolor,fancyhdr,titlesec}
\usepackage{hyperref}
\usepackage{needspace}
\definecolor{azul}{RGB}{35,91,133}
\definecolor{rojo}{RGB}{162,39,50}
\definecolor{gris}{RGB}{70,82,92}
\hypersetup{colorlinks=true,linkcolor=azul,urlcolor=azul,pdfauthor={Aquiles Taylor Ramos Yapo},pdftitle={Minicentral hidroeléctrica Limbani - Anteproyecto R02}}
\urlstyle{same}
\setlength{\parindent}{0pt}
\setlength{\parskip}{4pt}
\setlength{\tabcolsep}{4pt}
\renewcommand{\arraystretch}{1.18}
\setlist{nosep,leftmargin=1.5em}
\setcounter{tocdepth}{1}
\setcounter{secnumdepth}{2}
\titleformat{\chapter}[display]{\normalfont\bfseries\color{azul}}{\large CAPÍTULO \thechapter}{6pt}{\LARGE}
\titlespacing*{\chapter}{0pt}{0pt}{14pt}
\titleformat{\section}{\large\bfseries\color{rojo}}{\thesection}{.7em}{}
\titleformat{\subsection}{\normalsize\bfseries\color{azul}}{\thesubsection}{.7em}{}
\fancyhf{}
\fancyhead[L]{\small\bfseries MINICENTRAL HIDROELÉCTRICA LIMBANI\\\footnotesize Centrales Hidroeléctricas | Primera unidad}
\fancyhead[R]{\footnotesize UNA PUNO\\R02 | 13.10.2026}
\fancyfoot[L]{\scriptsize Aquiles Taylor Ramos Yapo | 227368 | GRUPO A}
\fancyfoot[R]{\scriptsize Anteproyecto académico | \thepage}
\renewcommand{\headrulewidth}{.5pt}
\renewcommand{\footrulewidth}{.3pt}
\pagestyle{fancy}
\fancypagestyle{plain}{\fancyhf{}\fancyhead[L]{\small\bfseries MINICENTRAL HIDROELÉCTRICA LIMBANI\\\footnotesize Centrales Hidroeléctricas | Primera unidad}\fancyhead[R]{\footnotesize UNA PUNO\\R02 | 13.10.2026}\fancyfoot[L]{\scriptsize Aquiles Taylor Ramos Yapo | 227368 | GRUPO A}\fancyfoot[R]{\scriptsize Anteproyecto académico | \thepage}}
\newcommand{\estado}[1]{\par\noindent\colorbox{azul!8}{\parbox{\dimexpr\linewidth-2\fboxsep}{\small\textbf{Estado de ingeniería.} #1}}\par}
\newcommand{\fuente}[1]{{\small\color{gris}\textit{Fuente / trazabilidad: #1}}\par}
\newcolumntype{L}[1]{>{\raggedright\arraybackslash}p{#1}}
\emergencystretch=2em
\input{generado/valores}

\begin{document}
\hypersetup{pageanchor=false}
\begin{titlepage}
\thispagestyle{empty}
\centering
{\large\bfseries UNIVERSIDAD NACIONAL DEL ALTIPLANO\par}
\vspace{3mm}
{\small FACULTAD DE INGENIERÍA MECÁNICA ELÉCTRICA, ELECTRÓNICA Y SISTEMAS\par}
{\small ESCUELA PROFESIONAL DE INGENIERÍA MECÁNICA ELÉCTRICA\par}
\vspace{8mm}
\includegraphics[height=27mm]{figuras/escudo-unap.png}\par
\vspace{8mm}
{\large\bfseries CENTRALES HIDROELÉCTRICAS\par}
\textcolor{azul}{\rule{\textwidth}{1pt}}\par
\vspace{5mm}
{\LARGE\bfseries\color{azul} MINICENTRAL HIDROELÉCTRICA LIMBANI\par}
\vspace{4mm}
{\large Objetivo de estudio en punto de conexión: 150 kW\par}
\vspace{5mm}
{\large\bfseries ANTEPROYECTO DE PRIMERA UNIDAD\par}
{\normalsize Memoria, base normativa, cálculos justificativos,\\especificaciones y esquemas conceptuales\par}
\vspace{9mm}
\begin{tabular}{L{4.1cm}L{9cm}}
\textbf{Alumno} & Aquiles Taylor Ramos Yapo\\
\textbf{Código de estudiante} & 227368\\
\textbf{Responsable de asignatura} & MEZA MAROCHO GREGORIO\\
\textbf{Sección} & GRUPO A\\
\textbf{Localización propuesta} & Limbani, Sandia, Puno\\
\textbf{Revisión} & R02 - 13 de octubre de 2026\\
\textbf{Nivel del estudio} & Anteproyecto académico con hipótesis explícitas\\
\end{tabular}
\vfill
\textcolor{azul}{\rule{\textwidth}{.5pt}}\par
{\small PUNO - PERÚ | 2026\par}
\end{titlepage}

\hypersetup{pageanchor=true}
\pagenumbering{roman}
\chapter*{Resumen del anteproyecto}
\addcontentsline{toc}{chapter}{Resumen del anteproyecto}
Se estudia una minicentral de pasada en Limbani, Sandia, Puno, para la primera unidad de Centrales Hidroeléctricas. La meta de estudio es exportar 150 kW en el punto de conexión común (PCC); su adopción como requisito del curso debe confirmarse con el docente. Se plantea operación en paralelo con red de 22.9 kV, conservando como hipótesis los niveles y el caudal del borrador recibido.

La alternativa prioritaria de estudio R02 usa $Q=0.420$ m$^3$/s, salto bruto 52 m, tubería de 220 m y diámetro interior 500 mm. Con Darcy-Colebrook y rendimientos supuestos, resulta salto neto \HnAlt\ m, potencia en eje de turbina \PtAlt\ kW y potencia de generador \PeAlt\ kW. El inventario operativo suma \PAuxNormal\ kW; se añaden pérdidas separadas del TSA, transformador principal y línea. El balance entrega \PExport\ kW en PCC, con margen de solo \MargenExport\ kW sobre la meta.

La opción de 450 mm entrega \PExportRef\ kW con el mismo inventario y caudal. Aumentar diámetro reduce pérdidas, pero no resuelve por sí solo la disponibilidad de agua ni la eficiencia real. La sensibilidad a un caudal 5\% menor hace caer la exportación por debajo de 150 kW. Una reserva de estudio del 5\% en potencia exigiría \QMargen\ l/s con 500 mm, superior al saldo no verificado de 420 l/s.

Se desarrollan conducción libre, reja, desarenado, alivio, transitorios, esfuerzos de tubería, apoyos, balance de carga, neutros, protecciones y metrado. Se comparan 450, 500 y 550 mm con las mismas fronteras energéticas. Las obras y equipos se especifican por desempeño y evidencia de aceptación. La ruta normativa enlaza cada requisito del CNE y legislación con el cálculo, plano, trámite o ensayo que lo puede demostrar.

\estado{Resultado: anteproyecto académico condicionado. No se ha acreditado caudal, topografía, derechos de agua, clasificación ambiental, conexión, material ni suelo. Ninguna alternativa está aprobada para construcción. Los cálculos de membrana y viga son cribados; transitorios reales, estabilidad local, cortocircuitos, puesta a tierra y energía anual requieren entradas adicionales.}

La revisión del borrador y de R01 se conserva en un informe separado. El cuerpo principal desarrolla la propuesta R02. Todas las cifras nuevas provienen de entradas identificadas y un constructor reproducible; las series hidrológicas y demanda externa permanecen vacías cuando no existe evidencia.
\begingroup\small\setlength{\parskip}{0pt}
\tableofcontents
\endgroup
\clearpage\pagenumbering{arabic}
\chapter{Memoria descriptiva y bases del proyecto}
\section{Objeto y alcance}
Desarrollar un anteproyecto académico de minicentral hidroeléctrica de pasada, evaluar la coherencia del borrador y establecer el camino verificable hacia un diseño sustentado. La potencia de 150 kW se conserva como objetivo de estudio procedente del borrador, pendiente de confirmar como requisito del docente. La ubicación en Limbani y el enlace propuesto con Electro Puno tienen el mismo carácter preliminar.

El alcance comprende captación, desarenado, conducción, cámara de carga, tubería forzada, grupo turbina-generador, servicios propios, transformación y enlace. Las obras se analizan como sistema: el agua disponible limita la energía; la carga y el punto de conexión condicionan potencia, regulación y protecciones; la topografía condiciona trazado, pérdidas y estabilidad.

\section{Documentos de partida y trazabilidad}
El archivo recibido es \texttt{main.pdf}, elaborado previamente con Gemini. Se conserva una copia inalterada en fuentes locales y se registra su huella digital. El expediente de vivienda de Aquiles aporta la organización por capítulos, tipografía sin serifas, escudo, encabezado institucional y tablas. Los cálculos residenciales no se transfieren a generación.

Los datos se clasifican en: \textbf{declarado en el borrador}, \textbf{hipótesis de esta revisión}, \textbf{calculado} y \textbf{pendiente}. Un valor declarado no equivale a medición. Las salidas numéricas muestran más decimales para auditoría; la precisión física sigue limitada por las entradas.

\begin{longtable}{L{4cm}L{5cm}L{5.4cm}}
\toprule Entrada & Información disponible & Tratamiento R02\\\midrule
Identidad académica & Alumno, código, docente y sección informados por el usuario & Confirmados para portada\\
Caudal y niveles & 0.380 m$^3$/s; cámara 3202 m y eje 3150 m & Hipótesis sin aforo ni topografía\\
Tubería & 175 m; interior 390 mm; DN 400 & Reconstrucción para auditoría, geometría inconsistente\\
Conexión & 22.9 kV; línea 1.2 km & Propuesta sin punto aprobado\\
Demanda de usuarios & No adjunta & No se inventan viviendas ni industrias beneficiarias\\
Normativa & CNE Utilización local y fichas parciales & Descarga CNE Suministro y fuentes legales; matriz por requisito\\
\bottomrule
\end{longtable}

\section{Ubicación y control geométrico}
La denominación política Limbani-Sandia-Puno se toma del borrador. El sector exacto, margen, predios, cuenca delimitada, sistema geodésico y servidumbres deben validarse. Las imágenes del borrador no tienen procedencia suficiente para demostrar el emplazamiento; se conservan en el original, sin usarlas como levantamiento.

Las coordenadas declaradas (Este, Norte, cota) están expresadas en metros, con UTM 19S declarado:
\[
\text{Cámara: }(425120,8441860,3202),\qquad
\text{Casa: }(425260,8441710,3150).
\]
Su distancia mínima es:
\begin{equation}
L_{\min}=\sqrt{140^2+150^2+52^2}=\MinLong\ \mathrm{m}>175\ \mathrm{m}.
\end{equation}
Una tubería real, con curvas y apoyos, no puede tener menos longitud que esa recta. La distancia horizontal es 205.18 m y la pendiente media geométrica aproximada es $\arctan(52/205.18)=14.22^\circ$. No corresponde afirmar simultáneamente 175 m y 18 grados para esos puntos.

Entre desarenador y cámara las coordenadas dan $\sqrt{210^2+200^2}=290.00$ m horizontales, por lo que tampoco respaldan un canal de 220 m. La longitud definitiva del canal queda pendiente. Los 220 m del escenario nuevo se refieren exclusivamente a la \textbf{tubería forzada}.

\section{Esquema funcional y modos de operación}
Se estudia captación con paso de caudal ecológico, desarenador con purga, canal a superficie libre, cámara con aliviadero, tubería, válvula de guarda, turbina de flujo cruzado candidata y restitución al cauce. La alternativa de bocatoma tirolesa debe compararse con una captación lateral considerando sedimentos, avenidas, mantenimiento y comportamiento ambiental.

\textbf{Modo principal propuesto:} paralelo a red, con potencia activa controlada por agua admitida y potencia reactiva por excitación según acuerdo de conexión. \textbf{Modo aislado:} alternativa futura, que exigiría demanda medida, control de frecuencia y tensión, reservas, rechazo de carga y eventualmente carga de lastre. No se presupone capacidad automática de operar en ambos modos.

La protección anti-isla debe desconectar el enlace ante pérdida de red conforme al estudio aprobado. Reiniciar y sincronizar exige comprobar condiciones de red, estado hidráulico y ausencia de fallas. La energía de maniobra y disparo debe permanecer disponible ante pérdida de servicios alternos.

\section{Criterios de aceptación de esta revisión}
Se acepta la coherencia algebraica y documental cuando las ecuaciones cierran, las unidades son consistentes, cada entrada tiene origen y las limitaciones están indicadas. Se difiere la aceptación física hasta disponer de mediciones y documentos. La aprobación académica corresponde al docente; un expediente para obra requeriría ingeniería, firmas, permisos y revisiones adicionales [N01, 011.F].

\chapter{Marco normativo y ruta de cumplimiento}
\section{Jerarquía y forma de aplicar las normas}
Las leyes y reglamentos peruanos establecen obligaciones. El CNE fija requisitos de seguridad eléctrica; no sustituye los estudios hidráulicos, geotécnicos ni la selección de maquinaria. Las normas IEC, IEEE, ASME y de producto deben emplearse dentro de su alcance y con la edición contractual acordada. Una lista bibliográfica no demuestra cumplimiento.

Esta revisión vincula regla, criterio, evidencia y estado. Se distinguen la \textbf{corrección de una referencia}, la \textbf{especificación de una exigencia} y la \textbf{verificación de su satisfacción}. La matriz del Anexo A conserva abiertos los requisitos que dependen de campo, fabricante o autoridad.

\section{Código Nacional de Electricidad: alcance concreto}
\subsection{Suministro y utilización}
El CNE Suministro 2011 se aprobó mediante R.M. 214-2011-MEM/DM [N01]. Sus reglas 011.A-C incluyen generación y sistemas interactivos con suministro público, y establecen obligatoriedad en Perú. La regla 011.B remite las instalaciones del usuario en edificaciones al CNE Utilización. La central y su enlace requieren Suministro; las instalaciones interiores de servicios propios requieren compatibilizar Utilización y las reglas de estación aplicables.

El CNE Utilización fue aprobado por R.M. 037-2006-MEM/DM [N02]. Se leyó junto a la modificación R.M. 175-2008-MEM/DM [N03], especialmente propagación de fuego, diferenciales y puesta a tierra. Los factores de demanda de vivienda no se aplican automáticamente a bombas, gobernador, cargador o máquinas auxiliares.

\subsection{Reglas que cambian el diseño}
\begin{itemize}
\item \textbf{CNE-S 130.A, B y E:} además del regulador, considerar disparo independiente por sobrevelocidad cuando pueda resultar perjudicial, parada manual accesible y circuitos de parada protegidos.
\item \textbf{CNE-S 132:} el neutro de generador requiere resolver la impedancia y las necesidades particulares del sistema. No se adopta neutro sólido por analogía con una vivienda.
\item \textbf{CNE-S 036.A-B:} diseñar para paso, toque y potenciales transferidos. La regla 362 trata subidas de tuberías en poste; no sustenta un límite general de 5 ohm.
\item \textbf{CNE-S 110-124 y 152:} recinto, ventilación, humedad, iluminación, salidas, partes móviles, partes vivas y ubicación del transformador son verificaciones de diseño.
\item \textbf{CNE-U 030-004 y 050-104:} la sección del cable depende de instalación, temperatura, agrupamiento y régimen continuo; no basta dividir kVA entre tensión.
\item \textbf{CNE-U 050-102:} el 2.5\% individual y 4\% combinado corresponden a circuitos de utilización. La regulación del enlace de generación exige además requisitos de la distribuidora.
\end{itemize}

\subsection{Línea aérea y altitud}
Las distancias de la línea se evalúan en la Sección 23 (reglas 230, 232, 234 y tablas correspondientes). Las secciones 24, 25, 26 y 27 cubren grado de construcción, cargas, resistencia y aislamiento. Las secciones 21 y 22 no sustituyen esas tablas. Un poste de 12 m no acredita por sí mismo distancia al suelo ni resistencia.

La corrección por altitud exige leer el alcance de cada regla. Por ejemplo, 232.C.1.b actúa sobre la distancia adicional definida en 232.C.1.a para tensiones que sobrepasan 33 kV; no autoriza aumentar uniformemente toda distancia de una línea de 22.9 kV en 28\%. La coordinación de aislamiento externo de equipos se estudia con altitud, contaminación, humedad, impulso y especificación del fabricante; distancia de fuga y distancia en aire son verificaciones diferentes.

\section{Concesiones y obligaciones del titular}
El artículo 3 de la Ley de Concesiones Eléctricas utiliza el umbral \emph{mayor de 500 kW} para generación hidráulica. La inferencia para el objetivo de 150 kW es que no activa por sí solo ese supuesto de concesión definitiva. El artículo 7 conserva el cumplimiento técnico, ambiental, patrimonial y la obligación de informar inicio y características de operación [N04 y contraste N13]. La línea y las servidumbres tienen análisis propio; no se declara que todo el proyecto esté libre de autorizaciones.

El alumno es autor académico, no se le asigna automáticamente calidad de concesionario, propietario del agua o titular de obra. Antes de tramitar debe definirse titular, predios, modalidad de suministro, conexión y autoridad competente.

\section{Agua, caudal ecológico y obras}
La Ley 29338 regula derechos de uso. Sus artículos 44 y 47 sustentan el derecho y la licencia; el 53 exige disponibilidad en cantidad, calidad y oportunidad, preservación de caudales ecológicos y derechos de terceros, entre otros requisitos. El artículo 74 aborda faja marginal; el 81 la evaluación ambiental hídrica y el 104 las obras de infraestructura hidráulica [N05]. La acreditación de disponibilidad no equivale a licencia de uso ni a autorización para ejecutar obras.

El Reglamento de la Ley, artículos 153-155, reconoce permanencia y variación estacional del caudal ecológico [N06]. La R.J. 267-2019-ANA, artículo 6, vincula metodología a la categoría ambiental y establece un tratamiento para pequeños proyectos/no comprendidos en el listado SEIA coordinado con la AAA [N07]. Por tanto, no corresponde imponer ``20\% de Q95'' como regla universal.

\subsection{Secuencia de permisos y evidencia}
Se descargó el D.S. 023-2014-MINAGRI desde SPIJ [N19]. Sus artículos modificados 79, 81, 84 y 85 distinguen acreditación, autorización de obras y licencia. La acreditación no habilita a ejecutar obras ni usar agua. Las obras exigen antes el instrumento ambiental y la autorización sectorial cuando corresponde. La licencia requiere obras verificadas. No se atribuyen a la central procedimientos simplificados de otros tipos de proyecto.

\begin{longtable}{L{3.4cm}L{6.4cm}L{4.6cm}}
\toprule Etapa & Evidencia que debe prepararse & Condición del estudio\\\midrule
Disponibilidad & Balance por periodo, calidad, derechos, captación y restitución & No hay serie ni acreditación\\
Ambiente & Tipo de obra, evaluación del tramo derivado, caudal ecológico y certificación & Clasificación condicionada\\
Obras & Planos, predios, servidumbres, ambiente y requisito sectorial aplicable & Sin autorización ni planos ejecutivos\\
Licencia / operación & Verificación de obras y régimen de uso aprobado & Sin licencia\\
\bottomrule
\end{longtable}
El texto consultado es una modificación identificada, no consolidación de todas las reformas posteriores. La R.J. 007-2015-ANA y sus anexos deben acompañar el expediente administrativo; su publicación de aprobación no reemplaza los formatos completos. Antes de tramitar se confirma TUPA, autoridad y versión, sin deducir plazos vigentes de una copia histórica.

\section{Evaluación ambiental y territorio}
El D.S. 014-2019-EM, artículos 7-11 y Anexo 1, distingue estudios y clasificación [N08-N09]. En su texto base, una central de pasada sin estructura de regulación, de hasta 20 MW, ubicada en el supuesto B y con línea asociada de hasta 20 km, tiene clasificación anticipada DIA. Esta es una \textbf{ruta condicional}; no se ha demostrado que Limbani cumpla todas las condiciones B ni que la captación no genere embalse. No corresponde asignar DIA solo porque el objetivo es 150 kW.

Se verificará superposición con ANP y amortiguamiento, ACR, ecosistemas frágiles, derechos colectivos, población y patrimonio. Un barraje que genere embalse puede cambiar la fila aplicable del Anexo 1. El enlace y las obras auxiliares se evalúan como parte integral del proyecto, evitando fraccionar impactos.

La consulta pública de enero de 2026 es un proyecto normativo [N14]. Además se identificó el D.S. 011-2026-EM que modifica el artículo 48 para actividades en curso [N15]. Se incorpora esta actualización sin convertir un procedimiento de adecuación de instalaciones existentes en autorización para una central nueva. La consolidación normativa completa a la fecha de ejecución sigue pendiente.

\section{Obras civiles, seguridad y normas de equipo}
El RNE [N10] requiere seleccionar las versiones vigentes para cargas, sismo, suelos, concreto, acero, seguridad y electricidad. La geotecnia y el perfil definen cimentaciones, anclajes y deslizamiento; no se aprueban cuatro bloques por simple conteo del borrador.

La Ley 29783, su reglamento y el RESESATE aprobado por R.M. 111-2013-MEM/DM forman la base de seguridad laboral [N16-N17]. El plan debe desarrollar bloqueo de energías eléctrica e hidráulica, trabajo en altura, izaje, excavación, inundación, soldadura y rescate. Las referencias internacionales se especifican con acceso autorizado; la ficha IEC consultada ya publica edición 60034-1:2026 [N11]. No se traslada como obligación una fórmula de desclasificación sin texto ni curva de fabricante.

\chapter{Recurso hídrico, demanda y energía}
\section{Estudio hidrológico necesario}
El borrador reporta área de cuenca, precipitación, caudal medio y Q95 sin estación, periodo, serie ni metodología de transferencia. La cercanía de Phara o Sandia no demuestra equivalencia hidrológica. Se requieren delimitación de cuenca en el punto de captación, registros SENAMHI/ANA con metadatos, control de faltantes, homogeneidad, aforos locales y método de regionalización calibrado.

Cada aforo debe registrar fecha, coordenadas, sección, nivel, equipo, calibración, procedimiento y error estimado. Una campaña en estiaje ayuda a contrastar el modelo, pero un aforo aislado no genera una curva de duración ni garantiza caudal todo el año. Deben revisarse derechos y captaciones aguas arriba y en el tramo derivado, retornos, avenida de diseño y transporte de sedimentos.

\section{Balance disponible y caudal de diseño}
Para cada intervalo $i$:
\begin{align}
Q_{disp,i}&=\max[0,Q_{rio,i}-Q_{eco,i}-Q_{terceros,i}-Q_{otros,i}],\\
Q_{turb,i}&=\min[Q_d,Q_{disp,i}].
\end{align}
El balance debe ser consistente con la sección de aforo: si el caudal observado ya refleja una extracción aguas arriba, no se descuenta nuevamente. Las purgas o pérdidas se incorporan donde realmente consumen o desvían agua. Cuando $Q_{turb,i}$ quede bajo el mínimo técnico del grupo, se analiza parada o funcionamiento parcial según fabricante.

Con los valores no comprobados del borrador, $0.520-0.100=0.420$ m$^3$/s. Este número coincide exactamente con el caudal de la alternativa y deja \textbf{margen nulo} para otros usos, incertidumbre y purgas. Por ello no respalda 150 kW firmes. Ni siquiera un Q95 validado representa 100\% del tiempo.

La curva de duración se obtendrá ordenando caudales de una serie definida, documentando periodo y posición de trazado (por ejemplo, $p=m/(n+1)$). Se debe diferenciar curva anual, mensual y de época seca. En esta revisión se genera una \emph{curva de potencia versus caudal}, no una falsa curva de duración del río.

\section{Demanda que justifica la central}
Hay dos problemas distintos. En paralelo a red, el objetivo puede ser exportar energía conforme a conexión y contrato; no implica atender un número inventado de viviendas. En modo aislado, se necesita inventario por usuario, demanda simultánea, perfil horario, cargas prioritarias, motores y crecimiento sustentado. Los 150 kW del borrador son una meta de oferta, no una demanda demostrada.

El levantamiento de cargas externas deberá contener usuario o grupo, cantidad, potencia eléctrica de entrada, factor de potencia, horas de uso, simultaneidad medida o justificada, arranque y prioridad. Para cada intervalo:
\[
P_{dem}(t)=\sum_j n_jP_{u,j}f_j(t),\qquad
Q_{dem}(t)=\sum_j P_j(t)\tan[\arccos(fp_j)].
\]
Se suman potencias reactivas algebraicamente; no se suman kVA indiscriminadamente ni se usa un único factor de demanda residencial para una central.

\section{Servicios propios propuestos}
El siguiente inventario es una \textbf{hipótesis de dimensionamiento}, con potencias eléctricas absorbidas. Para una bomba real cuya placa indique potencia mecánica se dividirá entre su eficiencia; no se duplicará esa corrección cuando el dato ya sea potencia eléctrica.
\small\input{generado/auxiliares}\normalsize
\small\input{generado/modos_auxiliares}\normalsize
En generación se consideran ocho luminarias interiores, cuatro exteriores, control, cargador, dos ventiladores, una bomba de drenaje y un gobernador activos: $P=\PAuxNormal$ kW y $Q=\QAuxNormal$ kvar. No se descuenta intermitencia de bombas sin un ciclo comprobado. La simultaneidad es una hipótesis operativa explícita, no un factor de demanda residencial. Mantenimiento suma $P=\PAuxMax$ kW y $S=\SAuxMax$ kVA con tomas e izaje, y se estudia con grupo parado e importación desde red.

Se conserva una bomba de reserva. Contar solo una exige enclavamiento y comprobar que la inundación de diseño no demanda ambas. La regla CNE-U 050-106(3) permite tratar cargas mutuamente excluyentes según la mayor, cuando esa exclusión está demostrada; 050-106(5) exige atender la máxima carga que pueda servirse para equipos cíclicos. No se reduce ampacidad simplemente por el consumo energético promedio.

El TSA candidato de 15 kVA alimenta auxiliares desde barra BT. Se suponen $P_{0,TSA}=0.08$ kW y $P_{cu,TSA}=0.25$ kW a placa. En generación sus pérdidas son \PerdTsa\ kW y se añaden al balance, sin confundirlas con carga de usuario. Sus potencias reactivas de magnetización y dispersión faltan en el modelo y deben incorporarse al flujo AC. La garantía en altura, corriente de arranque y caída transitoria de tensión pueden exigir otra potencia. Potencia continua suficiente no demuestra capacidad de arrancar el mayor motor.

\subsection{Arranque y transición entre modos}
Con red disponible y generador separado por 52G, el transformador principal puede alimentar la barra y el TSA si su esquema permite energización inversa. El control verifica agua, auxiliares, tensión, fase, velocidad y excitación antes de sincronizar. Se necesita enclavamiento que impida cierre fuera de sincronismo y recierre de red sobre una isla alimentada.

Al perder red, la CC debe mantener relés y disparos, pero una batería de control no demuestra capacidad para operar toda la unidad hidráulica del gobernador. Debe comprobarse acumulador hidráulico, retorno a posición segura o respaldo requerido para cerrar sin sobrepresión. El arranque autónomo queda sin demostrar: exige una fuente de servicios propios y una secuencia separada, además de control del sistema aislado.

\section{Energía anual y potencia firme}
La producción debe evaluarse con rendimientos variables y disponibilidad:
\begin{equation}
E_{net}=\sum_i\left[\rho g Q_{turb,i}H_{n,i}\eta_t\eta_m\eta_g/1000
-P_{aux,i}-P_{T,i}-P_{L,i}\right]\Delta t_i.
\end{equation}
La expresión se aplica en intervalos de operación; se modela por separado importación de auxiliares durante parada. Las horas fuera de servicio, vertimientos, límite del generador, restricciones de red y mínimos de turbina deben incorporarse. Si el despacho permite absorción de reactivos cambia la corriente y las pérdidas.

Como mera sensibilidad sobre el objetivo, $150\times8760\times FC$ da 657, 920 y 1117 MWh/año para $FC=0.50$, 0.70 y 0.85, redondeados. Ninguno es pronóstico del sitio. La energía firme, el factor de planta y la disponibilidad se mantienen sin valor validado hasta cerrar el recurso y operación.

\section{Campaña y datos de cierre}
\begin{enumerate}
\item Topografía georreferenciada y niveles de operación, incluyendo avenidas y restitución.
\item Series y aforos; inventario de derechos, usos y retornos; estudio ecológico según N05-N07.
\item Granulometría, concentración y mineralogía de sedimentos; temperatura del agua.
\item Registro horario de demanda o condiciones formales de exportación a red.
\item Inventario de auxiliares y garantías de desempeño del grupo para el rango de caudales.
\end{enumerate}

\chapter{Cálculos hidráulicos y selección del grupo}
\section{Frontera energética y constantes}
Se adopta agua con $\rho=1000$ kg/m$^3$, $g=9.81$ m/s$^2$ y $\nu=1.24\times10^{-6}$ m$^2$/s. Son constantes e hipótesis de cálculo, pendientes de temperatura del agua. La frontera se establece entre nivel de cámara y referencia de entrada/descarga de la turbina compatible con la garantía del fabricante:
\begin{equation}
H_b=Z_{cam}-Z_{ref}=3202-3150=52\ \mathrm{m}.
\end{equation}
Para la aproximación de flujo cruzado con descarga atmosférica se usa la referencia del eje como simplificación. La cota de restitución a 3148 m no se suma automáticamente a la energía utilizable: requiere configuración y recuperación efectiva en el equipo. Tampoco se usa el nivel de bocatoma como origen de la tubería a presión.

\section{Pérdidas y verificación independiente}
Para cada tramo y accesorios:
\begin{align}
A&=\pi D_i^2/4, & v&=Q/A, & Re&=vD_i/\nu,\\
\frac{1}{\sqrt f}&=-2\log_{10}\left(\frac{\epsilon/D_i}{3.7}+\frac{2.51}{Re\sqrt f}\right),\\
h_f&=f\frac{L}{D_i}\frac{v^2}{2g}, & h_{loc}&=\sum K\frac{v^2}{2g}, & H_n&=H_b-h_f-h_{loc}.
\end{align}
$f$ es el factor de Darcy, no el de Fanning. El motor resuelve Colebrook iterativamente; usa $64/Re$ en flujo laminar y rechaza el régimen de transición sin modelo específico. Se adopta rugosidad $\epsilon=0.045$ mm y $\sum K=1.0$ como sensibilidad de entrada, dos codos y válvula abierta. La distribución real de accesorios y envejecimiento pueden aumentar pérdidas.

El borrador calcula $h_f=3.51$ m con $f=0.0152$ y luego declara 2 m totales en otra sección. Aun conservando ese factor, el salto no sería 50 m desde una cámara con $H_b=52$ m. La solución automatizada usa las mismas ecuaciones en todas las tablas:
\small\input{generado/escenarios}\normalsize
\fuente{Entradas del borrador e hipótesis R01; salida calculos.json. El primer escenario reconstruye la aritmética, pero no resuelve la contradicción del trazado.}

\section{Alternativa para estudiar el objetivo}
Se estudia prioritariamente $Q=0.420$ m$^3$/s, $L=220$ m y $D_i=0.500$ m. Se comparan los demás diámetros con la misma cadena energética y auxiliares. No se ha seleccionado diámetro nominal comercial ni garantizado caudal. Resultan:
\begin{align}
v&=\VelAlt\ \mathrm{m/s}, & Re&=\ReAlt, & f&=\FAlt,\\
h_f&=\HfAlt\ \mathrm{m}, & h_{loc}&=\HlAlt\ \mathrm{m}, & H_n&=\HnAlt\ \mathrm{m}.
\end{align}
La cadena de conversión se mantiene explícita:
\begin{align}
\eta_{TG}&=0.84\times0.985\times0.925=0.765345,\\
P_h&=\rho gQH_n/1000=\PhAlt\ \mathrm{kW},\\
P_t&=\eta_tP_h=\PtAlt\ \mathrm{kW},\\
P_{acople}&=\eta_mP_t=\PmAlt\ \mathrm{kW},\\
P_G&=\eta_gP_{acople}=\PeAlt\ \mathrm{kW}.
\end{align}
La potencia nominal de una turbina se refiere al eje y la del generador a su salida eléctrica/capacidad térmica. Deben pedirse ambas, con el mismo punto de operación, y especificar qué pérdidas incluye cada garantía. La potencia exportada se calcula en el capítulo siguiente.

\subsection{Comparación de diámetros interiores}
\input{generado/diametros}
La tabla compara hidráulica y exportación a 0.420 m$^3$/s y 220 m. Usa los auxiliares del capítulo 3, por eso cambia respecto de R01. El diámetro de 500 mm mejora pérdidas, pero el margen neto sigue siendo \MargenExport\ kW. Se elegirá mediante costo presente de inversión, mantenimiento y energía perdida, con caudal variable y precios verificables. Los metros de tubería, espesor y pérdidas locales son hipótesis comunes, no diseños comerciales intercambiables.

\begin{figure}[htbp]\centering
\includegraphics[width=\linewidth]{figuras/sensibilidad.pdf}
\caption{Sensibilidad calculada. No representa disponibilidad hídrica ni eficiencia real a carga parcial.}
\end{figure}
La rugosidad, pérdidas locales y salto se someterán a sensibilidad. El margen del escenario objetivo es pequeño: una reducción de 1 m de salto neto a Q=0.420 m$^3$/s reduce la generación aproximadamente 3.15 kW. Una caída de tres puntos porcentuales de eficiencia de turbina (de 84\% a 81\%) reduce unos 5.55 kW cerca del punto ilustrado. Esa variación puede eliminar el margen de exportación.

\section{Captación, canal y desarenado}
\subsection{Captación y control de avenida}
La captación debe derivar caudal de diseño y permitir el caudal ecológico en el tramo derivado, con medición y protección frente a obstrucción. La avenida, socavación, flotantes y acceso de limpieza definen azud, compuertas y defensas. No se fija altura de barraje 1.20 m sin hidráulica fluvial y análisis ambiental.

Como ejemplo de rejilla de acceso (no diseño de azud tirolés), limitar velocidad normal a 0.60 m/s y admitir 40\% de área bloqueada exigiría:
\[
A_{bruta}\geq\frac{0.420}{0.60\,(0.60)\,(1-0.40)}=\AreaReja\ \mathrm{m^2}.
\]
Se distingue área bruta de área abierta. El 0.60 central es porosidad propuesta: $\phi=s/(s+t)$ para luz libre $s$ y ancho frontal de barra $t$. El otro 0.60 es velocidad de estudio en los huecos libres; 0.40 es fracción de bloqueo. R01 omitía la porosidad en su área bruta. Velocidad, porosidad y bloqueo no son valores normativos universales. Revisar pérdida limpia/sucia, separación según turbina y fauna, resistencia, limpieza y caudal máximo. La bocatoma tirolesa requiere modelo de admisión propio, no solo este cribado frontal.

\subsection{Canal a superficie libre}
Se estudia una sección rectangular independiente de la geometría no comprobada del borrador: ancho 0.80 m, profundidad total 0.85 m, pendiente 0.002 y $n=0.015$ de Manning, todos supuestos. Para $A=by$, $R=by/(b+2y)$:
\[
Q=\frac{1}{n}AR^{2/3}S^{1/2},\qquad Fr=\frac{Q/(by)}{\sqrt{gy}}.
\]
El tirante normal calculado es \Tirante\ m para 0.420 m$^3$/s, la velocidad \VCanal\ m/s y $Fr=\FrCanal$. A tirante límite de 0.60 m conduce \QCanal\ m$^3$/s y deja 0.25 m de borde libre. Esto verifica solo flujo uniforme: transiciones, curvas, remanso, pendiente disponible, disipación y arranque/parada pueden exigir mayor sección o borde libre.

El tirante crítico rectangular es $y_c=(Q^2/(gb^2))^{1/3}=\TiranteCritico$ m. El tirante normal excede $y_c$ y el flujo es subcrítico. El borde libre hasta altura de 0.85 m es \BordeLibre\ m, antes de oleaje, remanso y sedimentación. Manning supone flujo uniforme; entradas, curvas, transiciones y nivel de cámara requieren cálculo de perfil.

La pendiente implica caída de 0.002 m por metro. Entre desarenador y cámara, la distancia horizontal mínima declarada es 290 m y la diferencia de cotas es 1.5 m. Un canal de al menos esa longitud consumiría como mínimo 0.58 m por pendiente adoptada; el saldo sería como máximo 0.92 m antes de transiciones y variaciones de nivel. No se dispone de perfil para demostrar compatibilidad. Los 220 m del escenario son exclusivamente tubería forzada. No se fija concreto hasta contar con trazado, secciones y espesores.

\subsection{Desarenador: comprobación preliminar}
Se define como objetivo de estudio retirar partículas del orden de 0.20 mm, sujeto a la garantía de erosión de turbina. El tamaño no fija por sí solo velocidad de caída: se necesitan forma, densidad, temperatura y régimen. Se usa $w_s=0.025$ m/s únicamente como hipótesis de sensibilidad, sin aplicar Stokes fuera de su régimen.

Con factor hidráulico propuesto $F_h=1.5$, el área efectiva de sedimentación por cámara activa es:
\[
A_s\geq F_h\frac{Q}{w_s}=1.5\frac{0.420}{0.025}=25.2\ \mathrm{m^2}.
\]
Un ejemplo de cámara activa de 10 m por 2.6 m y tirante 1.5 m tiene área 26 m$^2$, volumen 39 m$^3$, velocidad longitudinal 0.108 m/s y residencia 92.9 s. La caída ideal desde 1.5 m tarda 60 s; la relación es 1.55. Dos cámaras de igual capacidad permiten una en servicio y otra en limpieza, si compuertas y canal de acceso lo permiten. No se ofrecen como dimensiones finales: revisar turbulencia, distribución, cortocircuito hidráulico, tolva, capacidad de sólidos y purga.

\subsection{Cámara de carga y alivio}
Una reserva de 25 m$^3$ del borrador equivale a $25/0.420=59.5$ s a caudal de diseño si todo el volumen fuese utilizable. El volumen operativo real excluye submergencia mínima, sedimentos y volumen muerto. La cámara no sustituye regulación de energía de larga duración.

El aliviadero debe evacuar la totalidad del caudal entrante cuando la turbina cierre, además de excedentes definidos por diseño. Con $Q=C_wbh^{3/2}$, $C_w=1.7$ m$^{1/2}$/s y carga $h=0.30$ m resulta $b\geq\AnchoVertedor$ m para 0.420 m$^3$/s. Coeficiente dimensional, aproximación, tipo de cresta, ahogamiento y descarga deben comprobarse. No equivale a avenida de diseño de la bocatoma. La entrada a tubería requiere submergencia y geometría que eviten vórtices y aire, sin imponer una profundidad arbitraria.

\section{Transitorios y tubería forzada}
\subsection{Celeridad y maniobra}
Se realiza un cribado elástico con módulo volumétrico del agua $K_w=2.2$ GPa, módulo del acero $E=200$ GPa y espesor candidato 6.35 mm. La restricción axial se simplifica a coeficiente unitario:
\[
a=\sqrt{\frac{K_w/\rho}{1+K_wD/(Ee)}}=\Celeridad\ \mathrm{m/s},\qquad
T_c=2L/a=\TCritico\ \mathrm{s}.
\]
Para cierre de 4 s, mayor que $T_c$, la aproximación de cierre lento tipo Michaud da:
\[
\Delta H_{lento}\simeq\frac{2Lv}{gT_s}=\DHlento\ \mathrm{m},\qquad
p_{max,cribado}=\PLento\ \mathrm{bar}.
\]
No es una garantía de Allievi completa. Para cierre abrupto ideal, Joukowsky daría $\Delta H=av/g=\DHrapido$ m y presión total aproximada \PRapido\ bar. El contraste demuestra por qué una falla de maniobra puede gobernar sobre el cierre normal; no representa una presión final de diseño ni implica que ese cierre ocurra realmente.

\estado{La simulación por características (MOC) queda pendiente de perfil, ley de apertura/área, inercia, regulador, aireación y condiciones de contorno. Debe obtener máximos y mínimos en toda la tubería, presión de vapor, separación de columna, sobrevelocidad y sensibilidad ante fallas.}

\subsection{Espesor, esfuerzos y apoyos}
El diámetro interior de 500 mm y espesor nominal 6.35 mm implican exterior 512.7 mm en la idealización; no se etiqueta como tubería comercial DN 500 sin catálogo. Deben definirse diámetro exterior comercial, espesor efectivo, material y soldadura.

La relación de membrana $\sigma_\theta\simeq pD/(2e_{ef})$ sirve para estimar esfuerzo circunferencial en pared delgada. La comparación requiere tensión admisible del código seleccionado y eficiencia de junta verificadas. No basta obtener un espesor pequeño: presión externa/vacío, ovalización, estabilidad, peso, flexión entre apoyos, térmica, transporte, fatiga y sismo pueden gobernar.

En un codo, el empuje por presión puede aproximarse por $F_p=2pA\sin(\theta/2)$, con fuerzas de cantidad de movimiento adicionales según el volumen de control. Los anclajes deben equilibrar esas acciones y comprobar deslizamiento, vuelco, presión admisible y estabilidad de ladera. La distancia de silletas depende de rigidez y esfuerzos, no se adopta 6 m sin cálculo.

\subsection{Espesor efectivo y contraste de maniobras}
Se toman tolerancia negativa 12.5\% y reserva de corrosión 1.00 mm como hipótesis de sensibilidad R02:
\[
e_{ef}=6.35(1-0.125)-1.00=\EspesorEfectivo\ \mathrm{mm}.
\]
No son descuentos atribuidos a una norma de producto. Se conserva diámetro exterior para obtener $D_{i,ef}=D_o-2e_{ef}$ y $D_m=D_o-e_{ef}$. La hidráulica principal usa el interior nuevo; vida final también requiere rugosidad y diámetro correspondientes.

Para pared delgada, con presión manométrica, $\sigma_\theta\simeq pD_m/(2e_{ef})$. En extremos cerrados la componente longitudinal ideal es $\sigma_z\simeq\sigma_\theta/2$. En 500 mm resulta \SigmaLenta\ MPa para cierre lento y \SigmaRapida\ MPa para Joukowsky abrupto. La aprobación exige material, temperatura, combinación, código y eficiencia de soldadura verificadas; no se aprueba 6.35 mm a partir de esta tabla.

\small\input{generado/mecanica}\normalsize
Las tres opciones usan cierre de 4 s, pared nominal 6.35 mm, módulos constantes y restricción axial simplificada. Para celeridad se usa pared nominal nueva; para membrana se aplica el descuento anterior. Una evaluación de vida final requiere actualizar celeridad y sección efectiva simultáneamente. El modelo no determina ocurrencia de un cierre instantáneo.

La presión rápida supera 16 bar en las tres opciones. Una especificación ``PN16'' no garantiza resistencia a ese evento. Válvulas, conexiones y protección contra sobrepresión deben justificarse con la maniobra posible y la presión de diseño del conjunto. Se requieren máximos/mínimos por MOC y clasificación de estados normal, rechazo y falla. La cámara de carga abierta no elimina la onda en la tubería; un bypass o válvula de alivio requiere dimensionamiento y prueba propios.

\subsection{Viga entre apoyos y deformación térmica}
Se estudia un tramo horizontal simplemente apoyado bajo peso propio efectivo y agua, sin refuerzos, pendientes, vibración, accesorios ni esfuerzos locales de silleta. Para sección anular y carga lineal $w$:
\[
I=\frac{\pi(D_o^4-D_{i,ef}^4)}{64},\quad
M_{max}=\frac{w\ell^2}{8},\quad
\sigma_b=\frac{M_{max}D_o}{2I},\quad
\delta_{max}=\frac{5w\ell^4}{384EI}.
\]
\small\input{generado/apoyos}\normalsize
Se reconoce la variación con vano; no se eligen 6 m para construcción. En inclinación se descompone el peso, se resuelven apoyos deslizantes y anclajes y se combinan esfuerzos. Se revisan ovalización, flexión local, vacío, estabilidad y sismo. USBR EM03 es referencia técnica histórica para los mecanismos de carga [N20]; no es reglamento peruano ni sustenta una tensión admisible moderna.

Para diferencia térmica de estudio 40 K, $\Delta L=\alpha L\Delta T=\Expansion$ mm a lo largo de 220 m. Un bloqueo axial completo produciría idealmente $E\alpha\Delta T=\TermicaBloqueada$ MPa. Estas acciones requieren juntas y apoyos coherentes. No se suman máximas de distintas condiciones sin definir la combinación física.

El cálculo conserva empuje de codo a 45 grados en resultados estructurados. En planta se requiere sumar vectores por presión y cantidad de movimiento, pesos, fricción de silletas y acciones sísmicas. Las cimentaciones necesitan capacidad portante, deslizamiento, vuelco y estabilidad global de ladera; la masa de un bloque no demuestra anclaje seguro.

\section{Turbina, velocidad y dimensiones preliminares}
La turbina de flujo cruzado es candidata por el rango de salto/caudal y la posibilidad de regulación por celdas; el fabricante N12 publica un rango que incluye el punto estudiado. Esto no acredita 84\% de eficiencia ni inmunidad a abrasión. Debe compararse con Francis y, si la selección de fabricante lo justifica, alternativas de impulso y velocidades distintas.

La velocidad síncrona satisface $n=120f/p$. A 60 Hz, 8 polos dan 900 rpm y 10 polos 720 rpm. Se estudia 900 rpm con acoplamiento directo, condicionado a tamaño y precio reales. La velocidad específica debe indicar convención, por ejemplo:
\[
n_s=\frac{n\sqrt{P_t\,[\mathrm{kW}]}}{H_n^{5/4}}.
\]
Es un índice convencional que cambia si se usa HP; no se comparan valores sin declarar unidades. No sirve para descartar una familia de turbina por un umbral aislado.

Con $C_v=0.97$ y $u/v_1=0.48$ como coeficientes académicos del borrador:
\[
v_1=C_v\sqrt{2gH_n},\quad u=0.48v_1,\quad D_r=\frac{60u}{\pi n}.
\]
En R02, $H_n=\HnAlt$ m produce $v_1=\VelChorro$ m/s y $D_r=\DiametroRodete$ m a 900 rpm. El índice específico de la convención anterior es \NsTurbina. Adoptar diámetro 0.300 m cambia la relación periférica y requiere rediseño de inyector. El ancho de 490 mm del borrador se puede contrastar por continuidad, sin repetir su coeficiente empírico 0.085:
\[
A_{chorro}=Q/v_1,\qquad s_{equiv}=\frac{Q}{Bv_1}=\EspesorChorro\ \mathrm{mm}\quad (B=490\ \mathrm{mm}).
\]
Es el espesor normal al chorro de una sección rectangular equivalente, no espesor de álabe ni dimensión de tobera aprobada. Arco de admisión, ángulos, obstrucción, distribución de agua y regulación por celdas requieren diseño propio. El par medio en eje sería
\[
T=P_t/\omega=\TorqueTurbina\ \mathrm{N\,m},\qquad \omega=2\pi(900/60)\ \mathrm{rad/s}.
\]
El acoplamiento y los ejes deben soportar además torsión transitoria, desalineación, fatiga, cargas radiales y velocidad de embalamiento. La velocidad síncrona no demuestra estabilidad del gobernador ante rechazo de carga.

La requisición debe solicitar potencia en eje de aproximadamente \PtAlt\ kW en el punto ilustrado, eficiencia garantizada por carga, velocidad de embalamiento y esfuerzos admisibles, inercia, aireación, sedimentología límite, disposición y pruebas. Número de álabes, espesor, materiales y soldaduras quedan a diseño del fabricante y revisión técnica.

\chapter{Diseño eléctrico, cargas y protecciones}
\section{Balance hasta el punto de conexión}
Se adopta 400 V entre fases, 60 Hz y $fp=0.85$ para el generador como hipótesis del borrador. Para el escenario alternativo:
\begin{align}
S_G&=P_G/fp=\SGen\ \mathrm{kVA},\\
Q_G&=P_G\tan(\arccos fp)=\QGen\ \mathrm{kvar},\\
I_G&=\frac{S_G}{\sqrt3\,0.4}=\IGen\ \mathrm{A}.
\end{align}
Para exactamente 150 kW y fp=0.85 se requieren 176.47 kVA y 254.71 A. La capacidad de placa de 200 kVA a 400 V representa 288.68 A; no significa 200 kW ni disponibilidad del recurso.

La exportación se evalúa desde una barra BT con auxiliares conectados a través de TSA. El modo de generación suma \PAuxNormal\ kW y \QAuxNormal\ kvar, más \PerdTsa\ kW de pérdidas TSA. El transformador principal candidato de 250 kVA tiene pérdidas hipotéticas $P_0=0.6$ kW y $P_{cu,n}=3.5$ kW, y el enlace mide 1.2 km como hipótesis. Se calcula $S_{carga}$ con $P_G-P_{aux}-P_{TSA}$ y $Q_G-Q_{aux}$; no se usa un fp auxiliar global elegido por separado. La reactiva de magnetización y dispersión de ambos transformadores está omitida y debe incorporarse al flujo de carga definitivo.
\begin{align}
P_T&=P_0+P_{cu,n}(S_{carga}/S_n)^2=\PerdTrafo\ \mathrm{kW},\\
P_{PCC}&=P_G-P_{aux}-P_{TSA}-P_T-P_L=\PExport\ \mathrm{kW}.
\end{align}
PCC denota punto de conexión común. El valor depende de fichas y del sentido de flujo de reactivos. No es una garantía comercial. Para asegurar 150 kW en PCC debe recalcularse recurso y potencia nominal con margen justificado; limitar $P_G$ a 150 kW produciría menos de 150 kW exportados.

Con 450 mm se requieren \QObjetivoRef\ l/s para exportar 150 kW, por encima del saldo no verificado de 420 l/s. Con 500 mm se requieren \QObjetivo\ l/s y el punto de estudio de 420 l/s exporta \PExport\ kW. El margen disponible es solo \MargenExport\ kW. Esta opción se lleva a estudio prioritario, pendiente de recurso, costo, trazado, garantía y condiciones de conexión.

\Needspace{19\baselineskip}
\subsection{Balance sin doble conteo}
\small\input{generado/balance_energia}\normalsize
La primera fila es la potencia bruta de agua. Las filas intermedias son pérdidas o consumo; la última es exportación. Su suma cierra la primera fila con residuo numérico menor que $10^{-9}$ kW. Un cierre aritmético verifica el modelo, no las entradas físicas.

\begin{figure}[htbp]\centering
\includegraphics[width=0.96\linewidth]{figuras/balance-energia.pdf}
\caption{De la potencia hidráulica bruta a la energía eléctrica exportable, en el punto R02.}
\end{figure}

\Needspace{12\baselineskip}
\subsection{Sensibilidad y criterio de reserva}
\small\input{generado/sensibilidad}\normalsize
Cada fila cambia una hipótesis aislada, manteniendo las demás. No son probabilidades ni una hidrología sintética. La rugosidad se contrasta con 0.15 mm; el rendimiento de turbina con 80\%. La eficiencia real también depende de carga y salto, por lo que las curvas constantes no sirven para estimar producción anual definitiva.

Como sensibilidad académica se estudia un margen adicional de 5\% sobre 150 kW. Con 500 mm exigiría \QMargen\ l/s, superior a 420 l/s. El 5\% no es mandato del CNE ni requisito confirmado del docente. La reserva debe justificarse por incertidumbre, envejecimiento y garantía. Si el contrato limita exportación a 150 kW, el controlador debe limitar caudal/potencia y la energía excedente no se cuenta como venta.

\section{Altitud y capacidad nominal}
El borrador escribe $k_h=1-(3150-1000)/10000$ y obtiene 0.892. Esa expresión vale \textbf{0.785}. La discrepancia aritmética es independiente de su validez física: no se ha encontrado sustento para tratarla como ley universal de IEC 60034-1. La edición 2026 está publicada según ficha IEC [N11]; la revisión requiere texto autorizado y garantía particular de la máquina.

Se propone cotizar un generador de \textbf{250 kVA} y un transformador de \textbf{250 kVA}, ambos con capacidad continua garantizada en la altitud y ambiente del sitio. Para el generador, la capacidad garantizada debe ser al menos \SGen\ kVA y \PeAlt\ kW en el punto estudiado. El factor mínimo se obtiene de $\SGen/250$ para placa 250 kVA, y de $\SGen/200$ para 200 kVA. No se convierten estos cocientes en factores de altitud de IEC. Se comprueban junto con temperatura, excitación, kW mecánicos, curva P-Q y pruebas en sitio.

Para comparar comercialmente se solicitarán garantías tanto de 200 como de 250 kVA en sitio, eficiencia, ventilación, elevación térmica, reactancias subtransitoria/transitoria y de secuencia, cortocircuito sostenido, constante de tiempo, inercia y soportabilidad de sobrevelocidad. Sobredimensionar kVA no crea potencia hidráulica.

\section{Transformador, grupo vectorial y neutros}
La denominación Dyn5 indica devanado de mayor tensión en delta (D) y de menor tensión en estrella con neutro accesible (yn); el número de reloj describe desplazamiento de fase conforme a la convención. Alimentarlo como elevador no convierte el delta de 22.9 kV en una estrella con neutro. La frase ``neutro de alta sólidamente aterrado'' del borrador se retira.

El grupo definitivo exige coordinación con Electro Puno: régimen de neutro de red, secuencia cero, puesta a tierra, sobretensiones temporales, protección de faltas y sincronismo. Se mantiene grupo vectorial \textbf{pendiente}. La rotación de fase, polaridades TC/TV y compensaciones de medición deberán corresponder al grupo real.

La regla CNE-S 132 conduce a estudiar impedancia del neutro del generador. El esquema conceptual no une automáticamente el neutro del generador con el del transformador y los auxiliares. Se propone transformador de aislamiento para servicios propios, para resolver su sistema de utilización sin anular la impedancia de puesta a tierra del grupo. El diseño de secuencia cero debe confirmar conexiones y protecciones. Los conductores de protección se integran equipotencialmente según el estudio; aislamiento funcional de circuitos no equivale a masas metálicas desconectadas de tierra.

\section{Conductores de potencia y auxiliares}
\subsection{Criterios de selección}
El diseño debe comprobar: corriente de servicio, ampacidad corregida, temperatura de terminales, operación continua, cortocircuito térmico, caída/elevación de tensión, instalación, resistencia mecánica, fuego y compatibilidad ambiental. La regla 050-104 distingue equipos marcados al 100\% o al 80\% y condiciones de tablas; no se aplica un factor 1.25 a todo circuito sin justificar.

En el tramo principal candidato, la placa de 250 kVA implica 360.84 A a 400 V. Si se diseña para esa capacidad, el cableado y barras deben corresponder a ella aunque el punto hidráulico solo demande \IGen\ A. Alternativamente se documentará una limitación permanente de potencia/corriente y se dimensionará coherentemente, sin usar ese límite como sustituto de protección.

\subsection{Lectura de carga continua del CNE Utilización}
El circuito de utilización se considera continuo salvo demostración de los intervalos de 050-104(3): para corriente hasta 225 A, no más de una hora total en dos; para más de 225 A, no más de tres horas en seis. En 050-104(1) la capacidad nominal del circuito es la menor entre protección y conductor.

Para marcado 100\%, 050-104(4) permite carga continua hasta 100\% con Tabla 2 o 85\% con Tabla 1; para marcado 80\%, 050-104(5) establece respectivamente 80\% y 70\%. Se verifican además las correcciones y la restricción subterránea de (6)-(7). La nota del Anexo B de 050-104 trata equipos sin marcado como 80\%. Estos requisitos de utilización se aplican dentro de su ámbito; para salida del generador se compatibilizan reglas de estación, fabricante y CNE-S, sin transferir un circuito doméstico como diseño de generación.

\begin{longtable}{L{4.6cm}rrrr}
\toprule Caso Cu 240 mm$^2$ & D8A (A) & D8B 100\% (A) & D8B 80\% (A) & Placa G (A)\\\midrule
Diagrama B4-1 detalle 4 & 446 & 442 & 442 & 360.84\\
Diagrama B4-1 detalle 7 & 390 & 386 & 386 & 360.84\\
\bottomrule
\end{longtable}
Las Tablas D8A y D8B se consultan juntas y con 030-004; no se reemplaza D8B simplemente por $0.8\times D8A$. Que un valor tabulado exceda 360.84 A no prueba capacidad en sitio: faltan temperaturas de conductor/terminal, suelo, agrupamiento, disposición y protección. El detalle 6 de D8A da 359 A, menor que la corriente de placa. No se selecciona un cable ejecutivo de 240 mm$^2$.

Para un ejemplo de utilización separado, Tabla 2, PVC 70 grados, método B2 y tres conductores cargados da 34 A para cobre 6 mm$^2$. Bajo hipótesis de 30 grados, sin agrupamiento y dispositivo marcado 80\%, un circuito nominal de 25 A tendría techo continuo 20 A, antes de otros factores. El ejemplo de 8 kW siguiente demanda 13.585 A y pasa ese contraste básico; no incluye arranque, protección de motores, neutro ni falla y no selecciona el alimentador del grupo.

\subsection{Ejemplo trazable de caída de tensión}
Para un alimentador trifásico auxiliar ilustrativo de 30 m, $P=8$ kW, fp=0.85, $U=400$ V y cobre de 6 mm$^2$, se usa resistividad a temperatura de servicio $\rho_e=0.0225$ ohm mm$^2$/m y reactancia $x=0.08$ ohm/km como supuestos:
\[
I=13.585\ \mathrm{A},\quad R=0.1125\ \Omega,\quad X=0.0024\ \Omega,
\]
\[
\Delta U=\sqrt3 I(R\cos\varphi+X\sin\varphi)\approx2.28\ \mathrm{V}=0.57\%.
\]
Se cumple el límite individual de 2.5\% en este ejemplo de utilización [N02, 050-102]; faltan el circuito derivado, neutro/desbalance, arranque y ampacidad para concluir. Si un derivado monofásico de 20 m, 10 A y 2.5 mm$^2$ usa los mismos supuestos resistivos, su caída aproximada resistiva es 3.6 V a 230 V (1.57\%). La suma ilustrativa es 2.14\%, menor que 4\%, siempre que ambos tramos y cargas correspondan al punto más alejado. No son recorridos medidos de la casa de máquinas.

Para esfuerzo térmico de falla puede estudiarse $S\ge I\sqrt{t}/k$ con material, aislamiento, temperaturas y $k$ documentados, o integrar $I^2t$ real de la protección. El cortocircuito decremental del generador requiere su curva; no basta extrapolar indefinidamente la corriente subtransitoria.

\section{Línea de enlace en 22.9 kV}
Se conserva 1.2 km como hipótesis. El conductor AAAC de 35 mm$^2$ del borrador es candidato, pendiente de ficha y cálculo mecánico. Se usan ilustrativamente $r=1.0$ ohm/km y $x=0.4$ ohm/km; no son valores certificados de ese conductor. Para longitud $\ell$:
\begin{align}
I_{MT}&=\frac{\sqrt{P^2+Q^2}}{\sqrt3 U}=\ILinea\ \mathrm{A},\\
P_L&=3I_{MT}^2r\ell/1000=\PerdLinea\ \mathrm{kW},\\
\delta U(\%)&\simeq\frac{r\ell P+x\ell Q}{10U^2}=\VariaLinea\ \%.
\end{align}
En la última ecuación P y Q se expresan en kW/kvar y U en kV. Para exportación, puede presentarse elevación de tensión en el extremo generador; el signo depende del flujo reactivo y de la convención. Se debe resolver flujo de carga con red, transformador, tomas, cargas locales y generación máxima/mínima.

La corriente baja no libera la línea de comprobar sección mínima aplicable, cortocircuito, flecha, tensión mecánica, vibración, rotura, vano de viento/peso, estructuras, cimentaciones, cruces y servidumbre. Se revisarán distancias del CNE-S para cada cruce a temperatura/flecha crítica, aislamiento y sobretensiones por la altitud. No se adopta descargador de 24 kV solo por tensión nominal de red: determinar MCOV/Uc, régimen de neutro, TOV, energía y coordinación.

\subsection{Consulta de distancias verticales hasta 23 kV}
La Tabla 232-1 del CNE-S, página impresa 126, contiene una columna de conductor de suministro expuesto de más de 750 V hasta 23 kV. Para una línea desnuda nominal 22.9 kV, y si se cumple la definición de tensión indicada en el encabezado, son referencias iniciales:
\begin{longtable}{L{10.7cm}r}
\toprule Superficie / cruce & Distancia (m)\\\midrule
Carreteras y avenidas sujetas a tráfico de camiones, fila 2.a & 7.0\\
Caminos, calles y otras áreas con camiones, fila 2.b & 6.5\\
Terrenos recorridos por vehículos, cultivos o pastos, fila 4 & 6.5\\
Áreas peatonales no transitables por vehículos, fila 5.a & 5.0\\
Calles y caminos en zonas rurales, fila 5.b & 6.5\\
\bottomrule
\end{longtable}
Se requieren lectura de notas y demás reglas, tensión máxima y régimen de puesta a tierra; no se usa automáticamente la tensión fase-tierra para cambiar de columna. Las distancias se comprueban con el conductor en condición crítica definida por 232.A, temperatura y flecha final, y con terreno levantado. Una altura de poste de 12 m debe descontar empotramiento, ubicación de herrajes y flecha; no demuestra por sí sola 7 m libres. Cruces de río, líneas, edificaciones y carreteras tienen verificaciones adicionales en 233/234 y notas de 232. El factor de altitud de 232.C.1.b no se aplica globalmente a estos valores. El expediente ejecutivo documentará cada vano y cruce con la regla y medición correspondientes.

\section{Cortocircuitos: órdenes de magnitud y límites}
Para un grupo candidato de 250 kVA, 400 V, $I_n=360.84$ A. Si $x_d''=0.15$ pu como hipótesis y tensión prefalla 1 pu:
\[
I''_{k,G}\approx I_n/x_d''=2.406\ \mathrm{kA}.
\]
Es contribución inicial de una fuente idealizada, no capacidad de interrupción necesaria. Si el transformador de 250 kVA tiene $u_k=4\%$ y la red MT fuese rígida, la contribución de red al defecto BT, despreciando línea, sería aproximadamente $I_n/0.04=9.021$ kA. La suma escalar daría un orden de magnitud de 11.4 kA, no un cálculo IEC de falla.

La selección exige modelo con potencia de cortocircuito de red, R/X, impedancias de secuencia, cableado, transformador, motores y generador; calcular máximos y mínimos, trifásico y fallas a tierra, corriente de pico, corriente de corte y evolución temporal. El interruptor debe cumplir Icu/Ics o capacidad aplicable, poder de cierre, servicio en sitio y selectividad. No se fija por ahora un interruptor ``10 kA'' ni fusible MT por el tamaño del portafusible.

\section{Protecciones y matriz causa-efecto}
Las funciones ANSI son identificación funcional, no ajustes automáticos. Se consideran 50/51 o variantes direccionales según flujos, falla a tierra acorde al neutro, 27/59, 81, 32, 40 y 25; evaluar 46, 49, 24, 87G/87T o alternativas según tamaño, sensibilidad y riesgos. La función 25 verifica sincronismo; no evita por sí sola una isla no intencional.

\begin{longtable}{L{3.6cm}L{5.5cm}L{5.3cm}}
\toprule Evento & Acción propuesta & Verificación pendiente\\\midrule
Pérdida de red / isla & Abrir enlace y bloquear reconexión; regular/cerrar agua según secuencia segura & Método anti-isla aceptado, tiempos y rechazo de carga\\
Sobrevelocidad & Disparo independiente; actuar agua y electricidad coordinadamente & Inercia, ley de cierre, presión y velocidad máxima\\
Falla interna del grupo & Abrir 52G, desexcitar, actuar válvula y bloquear rearranque & Sensibilidad, TC, zona y tiempo de despeje\\
Falla externa MT & Actuación selectiva en enlace, respaldo del grupo & Curvas, direccionalidad y recierre de distribuidora\\
Nivel bajo de cámara & Reducir carga o parar antes de aspirar aire & Sensores redundantes si procede, submergencia y temporización\\
Temperatura / vibración & Alarma y parada escalonada según criticidad & Límites garantizados y ubicación de sensores\\
Pérdida CA auxiliar & Mantener CC para relés, disparo y cierre seguro & Autonomía, cargador, vigilancia y pruebas\\
Parada de emergencia & Secuencia eléctrica/hidráulica definida, accesible y verificable & Riesgo de cierre brusco y energías residuales\\
\bottomrule
\end{longtable}
No se asignan umbrales finales de frecuencia ni tiempos de sobrecorriente sin estudio y coordinación. El ajuste de 60$\pm$1.5 Hz del borrador se retira como valor garantizado.

\section{Corriente continua y respaldo}
Se puede estudiar 24 Vcc o 48 Vcc según relés y actuadores. Como cálculo de capacidad, una carga crítica de 120 W durante 2 h, a 24 V, profundidad utilizable 0.8 y rendimiento 0.9 exige $120(2)/(24\times0.8\times0.9)=13.9$ Ah antes de envejecimiento, temperatura y consumos de disparo. Un factor de envejecimiento ilustrativo 1.25 elevaría a 17.4 Ah. No se selecciona batería hasta incluir corriente instantánea, número de maniobras, curva de descarga y ventilación [N01, sección 14].

\section{Puesta a tierra y equipotencialidad}
No se valida la malla 8 por 6 m ni seis electrodos del borrador. Se requieren mediciones Wenner en distintas separaciones y direcciones, estratigrafía y estacionalidad, corriente efectiva hacia tierra, reparto por neutros/pantallas, tiempo de despeje, capa superficial y geometría. El estudio evaluará resistencia, elevación de potencial y tensiones de paso, toque y transferidas [N01, 036.A-B].

El plano enlazará masas, bancada, carcasa, tubería, cercos, puertas, pantallas y descargadores; comprobará potenciales transferidos a obras alejadas. PE y neutro cumplen funciones distintas. Resistencia menor de 5 ohm no acredita seguridad. Guía y edición se acordarán contractualmente [N18].

\chapter{Especificaciones, montaje y pruebas}
\section{Especificación por desempeño}
Cada compra debe indicar condiciones de sitio, punto de operación, tolerancias, pruebas, documentación y garantías. La altitud declarada de 3150 m, clima y accesos necesitan confirmación. Se solicitará desempeño a condiciones reales, no solo placa a condiciones de referencia. Las capacidades siguientes son candidatas del escenario académico.

\begin{longtable}{L{3.2cm}L{6cm}L{5.2cm}}
\toprule Componente & Requisición preliminar & Evidencia de aceptación\\\midrule
Turbina & Flujo cruzado candidato, Q=0.420 m$^3$/s, Hn=\HnAlt\ m, potencia eje \PtAlt\ kW & Curvas garantizadas, planos, material, inercia, embalamiento y ensayos\\
Generador & Cotizar 200 y 250 kVA; verificar \SGen\ kVA y \PeAlt\ kW en sitio, 400 V, 60 Hz & Curva P-Q, pérdidas, térmica, reactancias, dinámica y certificados\\
Transformador & Candidato 250 kVA, 0.4/22.9 kV; grupo y aislamiento por conexión & Garantía en altura, pérdidas, impedancia, ensayo dieléctrico y rutina\\
Tubería & Interior modelo 500 mm; longitud modelo 220 m; espesor por diseño & Certificados de material, código aplicable, cálculo estructural y transitorio\\
Válvula de guarda & Presión y maniobra resultantes del estudio; bypass y posición segura & Curva Cv/apertura, ley de cierre, estanqueidad y energía de maniobra\\
Tableros BT/MT & Corriente, corte, arco y ambiente según estudio; acceso seguro & Verificación de diseño, ensayos, coordinación y planos de cableado\\
Servicios propios & TSA aislado candidato 15 kVA; CC para disparo/control & Inventario, arranques, selectividad, autonomía y neutros\\
Línea MT & Tensión 22.9 kV candidata; sección/estructuras por cálculo & Levantamiento, flechas, cruces, cimentaciones y aprobación de conexión\\
\bottomrule
\end{longtable}

\subsection{Balance térmico preliminar de casa de máquinas}
La pérdida activa del generador en el punto R02 es \PerdidaCalorGenerador\ kW. Si transmisión mecánica, auxiliares y TSA descargaran también calor dentro del recinto, la carga de estudio alcanzaría \CalorSala\ kW, sin transformador principal ni radiación solar. Se trata de una envolvente de cálculo, pendiente de ubicación y método de enfriamiento; no toda potencia auxiliar necesariamente se disipa en la sala.

El caudal de aire necesario por balance sensible sería $\dot V=\dot Q/(\rho_a c_p\Delta T)$. Como hipótesis explícita, $\rho_a=0.90$ kg/m$^3$, $c_p=1005$ J/(kg K) y salto térmico 10 K permiten evaluar el orden de magnitud. No se atribuye esa densidad a una medición ni se transforma 0.37 kW de ventilador en caudal de aire. Se pedirán curvas caudal-presión en altitud, pérdidas de rejillas/filtros, renovación, anticondensación y ventilación ante falla. El resultado térmico puede obligar a revisar la carga auxiliar, y esa revisión cambia la exportación.

\section{Materiales y detalles que no pueden omitirse}
\subsection{Agua y estructuras}
La captación debe permitir limpieza, purga y paso del caudal ecológico con medición verificable. Se dimensionan rejas y compuertas por carga diferencial máxima, obstrucción y corrosión. El desarenador necesita acceso y maniobra que eviten enviar una descarga concentrada de sólidos sin control ambiental.

Concreto, acero de refuerzo, recubrimiento, juntas, estanqueidad y cimentación se especificarán con diseño estructural y exposición. Resistencia de concreto de 210 kg/cm$^2$ del borrador no sustituye verificación de flexión, cortante, fisuración, sismo y durabilidad. En cámara y canal se incluirán drenajes, barandas, rescate y tapas cuando corresponda.

\subsection{Tubería y soldadura}
Se fijará código de diseño apropiado para la tubería forzada y límites de responsabilidad entre taller, montaje y fabricante. ASME B31.3 es un código de tuberías de proceso y no se declara automáticamente aplicable a todo el conducto hidroeléctrico; AWWA de producto y ASME IX de calificación tampoco reemplazan el cálculo estructural del sistema.

El expediente de fabricación debe incluir geometría por tramo, trazabilidad de coladas, tolerancias, procedimientos WPS/PQR, calificación de soldadores, mapa de soldaduras, recubrimiento interior/exterior, reparación y plan de inspección. El alcance END será el exigido por código, junta y criticidad, con aceptación definida antes de fabricar. Se retira el ``25\% de radiografía'' sin fundamento.

La prueba hidráulica debe considerar presión local a cada elevación, límites de todos los componentes, apoyos provisionales, purga de aire, temperatura, instrumentos calibrados y exclusión de personal. La presión y duración se fijarán mediante código y fabricante después del diseño; el producto $1.5\,p$ y 4 horas del borrador no constituyen protocolo aprobado. No se presuriza con aire como sustitución informal.

\subsection{Equipamiento eléctrico}
La envolvente se seleccionará por humedad, polvo, agua y ventilación. IP23 no se describirá como hermética ni cerrada estanca; equipos en un ambiente húmedo requieren verificar el grado adecuado y calefacción anticondensación si procede. Terminales, aislamiento y distancias se revisan en altura.

En utilización, aplicar N03: conductores no propagadores de llama y condiciones adicionales de fuego según uso; protección diferencial conforme a 020-132, incluyendo umbral no mayor de 30 mA donde corresponda a equipos accesibles a personas no calificadas. Seleccionar tipo y coordinación de diferencial según cargas; nunca reemplaza la puesta a tierra ni se instala indiscriminadamente en el circuito de generación de estación como si fuese un circuito doméstico.

La ubicación del transformador en aceite requiere evaluar separación de incendio, contención y drenaje que evite contaminación, acceso y señalización [N01, 152 y N08]. Se requieren TC/TV con clase, relación, saturación, carga secundaria, protección y puesta a tierra de circuitos secundarios definidas por el estudio.

\section{Montaje y control de calidad}
La liberación de montaje requiere planos aprobados, recepción de materiales y topografía de control. Se comprueban replanteo, dimensiones de cimentación y pernos, nivelación, grout, alineación y expansión térmica. Las tolerancias de alineación y vibración se tomarán del fabricante y norma aplicable al conjunto; se eliminan límites de 0.03 mm o 1.8 mm/s presentados como universales.

El plan de izaje define masas, centro de gravedad, radios, suelo, puntos certificados, accesorios y zona de exclusión. La tubería se monta con secuencia compatible con anclajes, no forzando piezas a alinearse. Se registra torque, soldadura, pruebas, recubrimiento y condición final de juntas de expansión.

El bloqueo debe aislar fuentes de red, generador, baterías, agua a presión y energía mecánica residual. Antes de intervención se comprueba ausencia de tensión, puesta a tierra temporal cuando proceda, imposibilidad de giro y despresurización. La puesta en servicio exige responsables identificados y control de accesos.

\section{Plan de ensayos y puntos de detención}
\begin{longtable}{L{3.2cm}L{6.2cm}L{5cm}}
\toprule Etapa & Verificación & Criterio / punto de detención\\\midrule
Recepción civil & Niveles, geotecnia, resistencia, pernos, drenaje & No instalar grupo sin liberación de cimentación\\
Recepción de tubería & Material, dimensiones, END, juntas y recubrimiento & No llenar sin soportes, purga y plan aprobados\\
Prueba hidráulica & Estanqueidad y estabilidad con instrumentos calibrados & Presión/duración definidas por diseño; acta por tramo\\
Eléctrica sin tensión & Continuidad, aislamiento, polaridad, fases y torque & Desconectar electrónica sensible según fabricante; límites de ensayo específicos\\
Protecciones & Inyección secundaria/primaria según alcance, lógica, disparos e interbloqueos & Valores aprobados; probar cadena completa hasta interruptor/actuador\\
Tierra & Resistividad, continuidad, medición y modelo paso/toque & No energizar si el estudio o ensayo no satisface criterios\\
Rotación y agua & Fugas, vibración, velocidad, temperatura y control & Límites del fabricante; confirmar presión mínima y máxima\\
Sincronización & Secuencia, desfase, tensión, frecuencia y permisos & Presencia/aceptación de distribuidora según procedimiento\\
Carga progresiva & Potencia, pérdidas, calentamiento y reactivos & Puntos contractuales y estabilización térmica definida\\
Rechazo de carga & Coordinación agua/eléctrica, sobrevelocidad y presión & Solo después de simulación y protocolo autorizado; registrar transitorios\\
Recepción & Manuales, planos conforme a obra, pruebas y capacitación & Sin observaciones críticas abiertas\\
\bottomrule
\end{longtable}
IEC 60041/60193 y normas de vibración se seleccionarán según ensayo de campo/modelo, tamaño y alcance contractual. No se obliga un ensayo de modelo universal para esta minicentral ni se confunde norma de prueba con garantía de rendimiento.

\section{Operación, mantenimiento y ambiente}
Se propone registrar niveles y caudales, energía bruta/neta, consumo propio, vertimiento, temperaturas, vibración, disparos y sedimentación. El mantenimiento se basa en horas, condiciones y fabricante: limpieza de rejas/desarenador, inspección de corrosión, apriete, aislamiento, aceite, batería y calibración de protección.

La operación debe preservar caudal ecológico y derechos, manejar purgas, aceites, residuos y ruido, y mantener respuesta a avenida, derrame, incendio y pérdida de red. Los intervalos de mantenimiento y stock crítico se definen con fabricante y accesibilidad; no se aseguran 100\% de disponibilidad ni costos mínimos por tipo de turbina.

\chapter{Metrados, evaluación económica y programa}
\section{Criterio de metrado}
El metrado debe proceder de geometría y especificaciones vigentes, no de porcentajes globales. Las cantidades siguientes son una base académica del escenario, no una orden de compra. Cada partida debe registrar código, unidad, cálculo, plano de origen, estado y revisión.

\small\input{generado/metrados}\normalsize
La masa de tubería se obtiene con $m=\rho_s\pi(D_o^2-D_i^2)L/4$, usando $D_o=D_i+2e=0.5127$ m. Debe recalcularse con despiece, tolerancia, espesores por tramo y accesorios; no reemplaza peso de catálogo.

\section{Análisis de precios y presupuesto}
El presupuesto de S/ 1,150,000 del borrador se conserva exclusivamente como antecedente no verificado. Sus proporciones no constituyen ACU. El ratio S/ 7,666.67 por kW reproduce la división entre 150 kW, pero no demuestra viabilidad ni equivalencia en dólares sin fecha y tipo de cambio.

Para cada partida $j$:
\[
CU_j=\sum_i q_{ij}p_{ij}+MO_j+EQ_j+TRAN_j,\qquad CD=\sum_j M_jCU_j.
\]
Se documentará si transporte está en precio de suministro o en partida separada para evitar doble conteo. Mano de obra depende de cuadrilla y rendimiento; equipo de horas y tarifa. Los precios deben registrar proveedor, fecha, moneda, impuestos incluidos/excluidos, lugar de entrega, vigencia, garantía y alcance. La oferta de turbina puede incluir o no gobernador, bancada, válvula, generador y pruebas.

Sobre costo directo se tratarán por separado ingeniería, supervisión, gestión ambiental/social, predios, conexión, gastos generales, contingencias justificadas e impuestos según el caso. No se completa artificialmente un total cuando faltan cantidades o cotizaciones. En esta revisión las columnas de precio y costo definitivo quedan pendientes en el archivo de metrado.

\section{Evaluación energética y económica}
Para evaluar inversión se requieren energía neta, venta o ahorro real, indisponibilidad, curtailment/restricciones, costos de operación, reposiciones, seguros, obligaciones hídricas, conexión y horizonte. Un costo por kW bajo no garantiza rentabilidad:
\[
VAN=-I_0+\sum_{t=1}^{N}\frac{Ingresos_t-OPEX_t-Reposiciones_t}{(1+r)^t}
+\frac{Valor_{residual}}{(1+r)^N}.
\]
No se asignan tarifa, tasa de descuento ni ingresos sin fuente. Se recomienda comparar al menos caudal bajo/base/alto, inversión, precio de energía y períodos fuera de servicio. El diámetro económico se evalúa con ese mismo horizonte y curva real de caudales.

\subsection{Decisión económica entre 450 y 500 mm}
El incremento puntual de exportación es \GananciaDiametro\ kW y el acero de envolvente aumenta \DeltaMasa\ kg con las mismas longitud y pared nominal. El sobrecosto real también incluye diámetro comercial, válvula, curvas, juntas, silletas, transporte e instalación; no se estima con kg adicionales solamente.

Para mostrar cómo comparar sin inventar tarifa se adoptan solo como sensibilidad 20 años, tasa real 10\% y 4000/6000/8000 horas equivalentes al punto estudiado. El factor de anualidad es $F=(1-(1+r)^{-N})/r=\FactorAnualidad$. Ninguna de esas horas es pronóstico de Limbani.
\small\input{generado/economia_diametro}\normalsize
Si ambas alternativas sostuvieran el punto el mismo número de horas, $\Delta E=\Delta P\,h$. Con límite de exportación de 150 kW se usa $\Delta P=\min(150,P_{500})-\min(150,P_{450})$, por lo que la energía marginal disminuye. Para una curva real, se integra por intervalo y carga, no se multiplica todo el año por el caudal de diseño.

Con tarifa neta verificable $\pi_e$ en S/kWh, el máximo sobrecosto justificable, antes de impuestos y con operación anual uniforme, sería:
\[
\Delta C_{max}=F(\Delta E\,\pi_e-\Delta OPEX)+\frac{\Delta V_{res}}{(1+r)^N}.
\]
La tabla permite evaluar ese umbral cuando existan tarifa y oferta comparables. No se produce VAN ni rentabilidad del proyecto mientras falten inversión y energía neta sustentadas; aumentar capacidad sin poder venderla puede empeorar la economía.

\section{Programa académico y ruta de proyecto}
El siguiente programa expresa dependencias propuestas, sin inventar una fecha de entrega del curso:
\begin{longtable}{L{1cm}L{4.1cm}L{4.4cm}L{4.9cm}}
\toprule Hito & Producto & Dependencia & Condición de cierre\\\midrule
H1 & Bases y requisitos docentes & Rúbrica y modo de operación & Potencia/alcance confirmados\\
H2 & Recurso, topografía y demanda & Acceso a datos y campo & Datos con incertidumbre y trazabilidad\\
H3 & Alternativas y viabilidad técnica & H2; consulta de conexión & Balance de agua/energía y opción seleccionada\\
H4 & Normativa y ambiente & Ubicación y tipo de obras & Ruta administrativa y clasificación confirmadas\\
H5 & Diseño detallado y ofertas & H3-H4; fabricantes & Transitorios, estructuras, cables, PAT y protecciones\\
H6 & Metrado, presupuesto y planos & H5 & Consistencia entre modelo y documentación\\
H7 & Revisión académica/técnica & Expediente integrado & Resolver observaciones críticas\\
H8 & Obra y puesta en servicio & Permisos y aprobación ejecutiva & Protocolos, ensayos y recepción\\
\bottomrule
\end{longtable}
Los ocho meses del borrador no se ratifican: permisos, campaña hidrológica, fabricación y restricciones de estiaje pueden controlar el plazo. Un cronograma de ejecución se formulará con rendimientos, suministros, ruta crítica y riesgos del sitio; no se confunde con las semanas disponibles para elaborar el trabajo académico.

\chapter{Esquemas, conclusiones y cierre de revisión}
\section{Esquemas producidos}
\begin{figure}[htbp]\centering
\includegraphics[width=\linewidth]{figuras/esquema-hidraulico.pdf}
\caption{H-00. Secuencia de aprovechamiento y responsabilidades del tramo derivado.}
\end{figure}
La captación, paso ecológico, purgas y restitución son partes de un mismo aprovechamiento. Las descargas requieren diseño de erosión, calidad y seguridad; no se da por inocua una purga porque vuelve al río.

\begin{figure}[htbp]\centering
\includegraphics[width=\linewidth]{figuras/perfil.pdf}
\caption{H-01. Perfil energético idealizado entre cámara y referencia de turbina. La línea no es levantamiento del terreno.}
\end{figure}
La línea de energía parte del nivel de cámara y desciende con pérdidas modeladas. La línea piezométrica en el interior de la tubería se obtiene restando $v^2/(2g)$ a la energía; la presión local depende además de la cota del eje. No se infiere seguridad frente a vacío a partir del perfil recto simplificado. El máximo transitorio no está dibujado como garantía.

\begin{figure}[htbp]\centering
\includegraphics[width=\linewidth]{figuras/unifilar-funcional.pdf}
\caption{E-01. Unifilar conceptual con barra BT y derivación a servicios propios. Sin ajustes ni secciones aprobadas.}
\end{figure}
E-01 permite revisar responsabilidades y secuencia del sistema. El unifilar definitivo deberá incluir número de polos, conductores, secciones, barras, tierra, neutros, TC/TV, protecciones, relaciones, enclavamientos, grupo vectorial y límites de propiedad. La simbología de construcción se normalizará con las referencias DGE y convenciones del proyecto.

\clearpage
\section{Planos que faltan para un expediente de detalle}
\begin{longtable}{L{1.5cm}L{5.7cm}L{7.2cm}}
\toprule Código & Plano requerido & Entrada que falta\\\midrule
PL-01 & Ubicación, predios y esquema general georreferenciado & Levantamiento y verificación de derechos/afectaciones\\
PL-02 & Perfil hidráulico y topográfico & Eje real, niveles, secciones y restitución\\
PL-03 & Captación y desarenador & Avenidas, sedimentos, geotecnia y diseño estructural\\
PL-04 & Cámara, rejas y aliviadero & Niveles de maniobra, vórtices, excedentes y acceso\\
PL-05 & Tubería y anclajes & Perfil, transitorio, código mecánico y suelos\\
PL-06 & Casa de máquinas y cimentación & Plano de fabricante, izaje, sismo y ventilación\\
PL-07 & Disposición del grupo & Fabricante, cargas dinámicas y mantenimiento\\
PL-08 & Unifilar y lógica de control & Estudio de conexión, fallas y protecciones\\
PL-09 & PAT, neutros y equipotencialidad & Resistividad, corrientes y tiempos de falla\\
PL-10 & Línea MT, perfiles y estructuras & Trazo, cruces, vanos, servidumbre y aislamiento\\
PL-11 & Auxiliares, iluminación y emergencia & Planta y cuadro de cargas validado\\
\bottomrule
\end{longtable}
Las escalas de estos planos se elegirán al disponer de geometría. No se anuncia que se adjuntan planos de detalle inexistentes.

\section{Conclusiones sustentadas}
\begin{enumerate}
\item El borrador no cierra su geometría y energía. Con su caudal, salto neto de 50 m y rendimientos, la salida es \POriginalCincuenta\ kW. Al recalcular su tubería interior de 390 mm y 175 m resulta \PBase\ kW, manteniendo aún una longitud incompatible con las coordenadas.
\item La alternativa prioritaria de 0.420 m$^3$/s, 220 m, diámetro interior 500 mm y salto bruto 52 m produce condicionalmente \PeAlt\ kW en bornes y \PExport\ kW en PCC. Su margen es \MargenExport\ kW. Con 450 mm, la misma carga auxiliar reduce la exportación a \PExportRef\ kW.
\item El caudal mínimo para 150 kW PCC sería \QObjetivo\ l/s con 500 mm bajo estas hipótesis. Una reducción de caudal del 5\% deja la exportación por debajo de la meta. No se acredita potencia firme ni garantía a partir del punto de diseño.
\item Los modos auxiliares explican \PAuxNormal\ kW en generación y \PAuxMax\ kW en mantenimiento; TSA, transformador principal y línea se contabilizan separados. La revisión térmica y de arranque puede exigir cambiar esos valores y recalcular exportación.
\item Las presiones de cribado son \PLento\ bar con cierre de 4 s y \PRapido\ bar para cierre abrupto ideal en 500 mm. Membrana, viga y térmica cuantifican acciones, pero no aprueban material, espesor, apoyos ni PN16 sin código y transitorio ejecutivo.
\item Los datos Q95 y caudal ecológico del borrador no sostienen disponibilidad anual. La nueva alternativa consume todo el saldo ilustrativo de 0.420 m$^3$/s, por lo que la hidrología y derechos deben cerrarse antes de seleccionar potencia.
\item La selección eléctrica requiere garantías en altura, régimen de neutro y puesta a tierra, estudio de conexión y cortocircuito. Se corrigen la descripción Dyn5, la fórmula de altitud y la referencia a Regla 362.
\item La matriz normativa permite revisar qué exige cada norma y qué evidencia falta. Ninguna declaración de cumplimiento integral se infiere de la inclusión de normas en bibliografía.
\item El proyecto queda organizado por capítulos con entradas canónicas y motor reproducible. Las verificaciones numéricas comprueban ecuaciones y coherencia; no reemplazan ensayos, levantamiento ni aprobación técnica.
\end{enumerate}

\section{Orden de cierre recomendado}
Primero se confirma con el docente si Limbani, 150 kW y conexión a red son requisitos o elecciones del alumno. Después se cierra el recurso y la geometría, se compara tecnología con fabricantes y se solicita información de conexión. Con estos datos se desarrollan transitorios, estructuras, protecciones y tierra; recién entonces se fijan metrados y costos definitivos.

La revisión normativa debe incorporar textos consolidados, las reglas complementarias de electrificación rural si la clasificación del sistema las hace aplicables y exigencias específicas de distribuidora. Se deberá obtener la rúbrica del curso para determinar extensión y profundidad final sin alterar la trazabilidad del estudio.

\section{Reproducción y control documental}
El punto de entrada es \texttt{proyecto.yaml}. El constructor lee los archivos JSON de parámetros, normativa y observaciones, calcula resultados y genera tablas y figuras en el directorio de construcción del proyecto. Las fuentes editables del informe están en \texttt{expediente/capitulos/}.

La ejecución produce \texttt{calculos.json}, comparación CSV, observaciones Markdown, huellas SHA-256 y los dos PDF. Se prueban balance energético, Colebrook, geometría imposible, entradas inválidas, maniobra rápida, flujo uniforme y cálculo trifásico. Los archivos técnicos permanecen como revisión R02 hasta validación del docente y de los datos del sitio.

\appendix
\chapter{Matriz de requisitos y evidencias}
\begingroup\footnotesize\renewcommand{\arraystretch}{1.08}
\input{generado/matriz}
\endgroup
\chapter{Fuentes y alcance de lectura}
Consulta: 13 de octubre de 2026. Los identificadores N01-N20 se utilizan en los capítulos y en la matriz. Los textos públicos descargados se conservan con huella SHA-256. Las fichas comerciales y las vistas previas de normas internacionales no se presentan como textos normativos completos.
\small\input{generado/fuentes}\normalsize
\end{document}
